\section{第10章\quad 重积分}
\subsection{二重积分的概念与性质}

\begin{exercise}{284}{积分方程与参数讨论}
设 $D=\{(x,y):0\leq y\leq x\leq1\}$，$\lambda\in\mathbb R$。讨论连续函数方程
\[
 f(x,y)=xy+\lambda\iint_D f(u,v)\,\dd u\dd v
\]
的解是否存在、是否唯一；存在时求出 $f(x,y)$，并给出 $\lambda=1$ 时的结果。
\end{exercise}

\begin{exercise}{285}{二重和的积分极限}
将下列极限分别表示为二重积分并计算，说明两者积分区域的差别：
\[
 \lim_{n\to\infty}\sum_{i=1}^n\sum_{j=1}^n
 \frac{n}{(n+i)(n^2+j^2)},\qquad
 \lim_{n\to\infty}\sum_{i=1}^n\sum_{j=1}^i
 \frac{n}{(n+i)(n^2+j^2)}.
\]
第二个极限只须化为一个定积分，无须求出初等形式。
\end{exercise}

\begin{exercise}{286}{黎曼和与误差估计}
计算
\[
 L=\lim_{n\to\infty}\frac{\pi}{2n^4}
 \sum_{i=1}^n\sum_{j=1}^n i^2\sin\frac{j\pi}{2n}.
\]
利用 $x^2$ 与 $\sin y$ 在相应区间上的单调性，说明该和式与 $L$ 之差为 $O(1/n)$。
\end{exercise}

\begin{exercise}{287}{积分比较与严格性}
设 $D$ 为顶点 $(1,0),(1,1),(0,1)$ 的三角形闭区域。比较
$\iint_D\ln(x+y)\,\dd\sigma$ 与 $\iint_D[\ln(x+y)]^2\,\dd\sigma$ 的大小，并证明所得不等式为严格不等式，说明边界上取等号是否影响结论。
\end{exercise}

\begin{exercise}{288}{幂函数比较与严格性}
设 $D=\{(x,y):(x-1)^2+(y-1)^2\leq2\}$，$u=(x+y)/4$。确定 $u$ 在 $D$ 上的取值范围，严格比较
\[
 I_1=\iint_Du\,\dd\sigma,\quad I_2=\iint_D\sqrt u\,\dd\sigma,\quad I_3=\iint_D\sqrt[3]u\,\dd\sigma.
\]
进一步证明 $J(p)=\iint_Du^p\,\dd\sigma$ 在 $p>0$ 时严格递减。
\end{exercise}

\begin{exercise}{289}{积分估值与分离变量}
设 $D=[0,\pi]\times[0,\pi]$。先用被积函数的最大、最小值估计
$I=\iint_D\sin^2x\sin^2y\,\dd\sigma$，再利用变量分离求其准确值；求 $\sin^2x\sin^2y$ 在 $D$ 上的平均值，并解释平均值与最大值不同的原因。
\end{exercise}

\begin{exercise}{290}{最优函数界与积分估值}
设 $D=\{(x,y):x^2+y^2\leq4\}$。求使 $m\leq x^2+4y^2+9\leq M$ 在 $D$ 上成立的最大 $m$ 和最小 $M$，据此估计二重积分；再计算其准确值，检验估计结果。
\end{exercise}

\begin{exercise}{291}{分片常值函数的积分}
设 $0<a<b<3$，$D=[0,3]\times[0,3]$，
\[
 f_{a,b}(x,y)=\begin{cases}1,&0\leq y\leq a,\\2,&a<y\leq b,\\3,&b<y\leq3.\end{cases}
\]
求 $\iint_D f_{a,b}\,\dd\sigma$，给出 $a=1,b=2$ 的结果；若只改变 $y=a$ 与 $y=b$ 两条线段上的有界函数值，积分是否改变？说明理由。
\end{exercise}

\begin{exercise}{292}{平均值与中值点}
设 $D=\{(x,y):0\leq y\leq x\leq1\}$。求 $f(x,y)=xy$ 在 $D$ 上的平均值，并求出 $D$ 内所有使 $f(\xi,\eta)$ 等于该平均值的点，说明其中是否存在区域内部的点。
\end{exercise}

\begin{exercise}{293}{缩小圆域上的平均值极限}
设 $D_r=\{(x,y):x^2+y^2\leq r^2\}$，$r>0$，
$A(r)=\frac1{\pi r^2}\iint_{D_r}e^{x^2-y^2}\cos(x+y)\,\dd x\dd y$。
求 $\lim_{r\to0^+}A(r)$；进一步利用二阶展开与圆域对称性，求 $\lim_{r\to0^+}[A(r)-1]/r^2$，并说明余项积分后仍可忽略的依据。
\end{exercise}

\begin{exercise}{294}{局部连续性与积分极限}
设 $f$ 在原点的某邻域内连续，$f(0,0)=1$；$g$ 在 $0$ 的某邻域内连续可微，且 $g(0)=0,g'(0)=3$。令 $D_r=\{x^2+y^2\leq r^2\}$。
求 $\lim_{r\to0^+}\frac{\iint_{D_r}f(x,y)\,\dd x\dd y}{g(r^2)}$，并用连续性的定义证明分子为 $\pi r^2+o(r^2)$，不能仅引用形式上的无穷小替换。
\end{exercise}

\begin{exercise}{295}{对称性使用的存在性条件}
有人利用奇偶性声称
$\iint_{[-1,1]\times[-2,2]}(x^2\tan y+\sin x^3)\,\dd\sigma=0$。
指出被积函数在区域内的奇点，并判断该式能否作为通常的有界函数二重积分成立。
再设 $0<h<\pi/2$，计算同一被积函数在 $[-1,1]\times[-h,h]$ 上的二重积分，写明对称性成立所需的条件。
\end{exercise}

\begin{exercise}{296}{变量互换与参数积分}
设 $D=\{(x,y):x^2+y^2\leq4,x\leq0,y\leq0\}$，$f$ 在 $[-2,0]$ 上连续且严格为正，$a,b\in\mathbb R$。
不用求 $f$ 的原函数，计算
$\iint_D\frac{af(x)+bf(y)}{f(x)+f(y)}\,\dd\sigma$。
给出积分关于 $a,b$ 的依赖关系，并说明 $a=b$ 时的结果与一般公式是否一致。
\end{exercise}

\begin{exercise}{297}{三角形区域上的对称分解}
设 $D=\{(x,y):-1\leq x\leq y\leq1\}$，$D_1=D\cap\{x\geq0,y\geq0\}$。
分别计算 $\iint_Dx^3y\,\dd\sigma$ 与 $\iint_D\cos x\sin y\,\dd\sigma$，从而求两者之和；证明后一个积分等于 $2\iint_{D_1}\cos x\sin y\,\dd\sigma$，指出所用的区域变换。
\end{exercise}

\begin{exercise}{298}{对称性与圆域积分}
设 $D=\{(x,y):x^2+y^2\leq1\}$。不先作二次积分，利用对称性求
\[
 \iint_D(x^2+2y)\,\dd x\dd y,
\]
并求 $\iint_D(ax^2+by^2+cxy+dy)\,\dd x\dd y$，其中 $a,b,c,d$ 为实常数。
\end{exercise}

\begin{exercise}{299}{保号性与中值定理的条件}
设 $D$ 为单位闭圆盘，$f$ 在 $D$ 上连续且非负。证明 $\iint_Df\,\dd\sigma=0$ 当且仅当 $f\equiv0$；若将连续条件去掉，仅要求有界可积，考察“原点取 $1$、其余点取 $0$”的函数，说明结论是否保持。
再令 $E$ 为两个互不相交的单位闭圆盘之并，在两圆盘上分别令 $g=0$ 与 $g=2$。判断能否找到 $P\in E$ 使 $g(P)$ 等于 $g$ 的积分平均值，并指出中值结论还需要的区域条件。
\end{exercise}

\begin{exercise}{300}{双矩积分方程的全部参数分支}
设 $D=[0,1]^2$，$\lambda\in\mathbb R$，连续函数 $f$ 满足
\[
 f(x,y)=x+y+\lambda\iint_D(xu+yv)f(u,v)\,\dd u\dd v.
\]
记 $A=\iint_Duf(u,v)\,\dd u\dd v$，$B=\iint_Dvf(u,v)\,\dd u\dd v$。建立关于 $A,B$ 的线性方程，完整讨论无解、唯一解及无穷多解的参数情形，写出存在时的全部 $f$，并代回验证。
\end{exercise}

\subsection{二重积分的计算方法}
\subsubsection{积分区域与交换积分次序}

\begin{exercise}{301}{交换积分次序与计算}
设 $f$ 在 $D=\{(x,y):0\leq x\leq4,\ x\leq y\leq2\sqrt{x}\}$ 上连续。
把 $\int_0^4\dd x\int_x^{2\sqrt{x}}f(x,y)\,\dd y$ 改写成先对 $x$ 积分的累次积分，并用所得表达式计算 $f(x,y)=y$ 时的积分。
\end{exercise}

\begin{exercise}{302}{半圆域与积分次序}
对连续函数 $f$，将
\[
 \int_0^1\dd y\int_{-\sqrt{1-y^2}}^{\sqrt{1-y^2}}f(x,y)\,\dd x
\]
交换积分次序。再取 $f(x,y)=x^2+y^2$，写出极坐标表达式并计算其值。
\end{exercise}

\begin{exercise}{303}{分块区域的合并}
设 $f$ 在下式涉及的闭区域上连续，将
\[
 \int_0^1\dd x\int_{-\sqrt{x}}^{\sqrt{x}}f(x,y)\,\dd y
 +\int_1^4\dd x\int_{x-2}^{\sqrt{x}}f(x,y)\,\dd y
\]
改写为一个先对 $x$ 积分的累次积分，并由此求该积分区域的面积。
\end{exercise}

\begin{exercise}{304}{分段边界与换序}
设 $f$ 连续，把
\[
 \int_0^1\dd y\int_0^{2y}f(x,y)\,\dd x
 +\int_1^3\dd y\int_0^{3-y}f(x,y)\,\dd x
\]
交换积分次序；取 $f(x,y)=y$，计算积分并求该区域的形心纵坐标。
\end{exercise}

\begin{exercise}{305}{三角函数边界的换序}
设 $f$ 在 $D=\{(x,y):0\leq x\leq\pi,-\sin(x/2)\leq y\leq\sin x\}$ 上连续。
将 $\iint_Df\,\dd x\dd y$ 改为先对 $x$ 积分的表达式，写明必须分段的纵坐标，并求 $D$ 的面积。
\end{exercise}

\begin{exercise}{306}{抛物线区域的分段换序}
设 $f$ 在 $D=\{(x,y):1\leq y\leq2,y\leq x\leq y^2\}$ 上连续。
将 $\int_1^2\dd y\int_y^{y^2}f(x,y)\,\dd x$ 交换积分次序，写明分段点；再取 $f(x,y)=x-y$ 计算，并从区域位置验证结果为正。
\end{exercise}

\begin{exercise}{307}{分块区域的重新组合}
设 $f$ 连续，记
\[
 I=\int_1^2\dd x\int_x^{x^2}f(x,y)\,\dd y
 +\int_2^8\dd x\int_x^8f(x,y)\,\dd y.
\]
分别用不等式描述两项对应的区域，将 $I$ 全部改为先对 $x$ 积分，并尽量合并具有相同积分限的部分；计算 $f\equiv1$ 时的值。
\end{exercise}

\begin{exercise}{308}{变限二重积分求导}
设 $f$ 在 $[1,+\infty)$ 上连续，$F(t)=\int_1^t\dd y\int_y^t f(x)\,\dd x$，$t\geq1$。
先交换积分次序，再求 $F'(t)$ 与 $F'(3)$；若 $f$ 处处为正，判断 $F$ 在 $(1,+\infty)$ 上的单调性并说明理由。
\end{exercise}

\subsubsection{直角坐标计算与分区}

\begin{exercise}{309}{矩形域与参数变化}
设 $a,b>0$，$D_{a,b}=[0,a]\times[0,b]$。求 $I(a,b)=\iint_{D_{a,b}}e^{x-y}\,\dd x\dd y$，给出 $I(1,1)$，并判断固定 $a$ 时 $I(a,b)$ 随 $b$ 的变化及其极限。
\end{exercise}

\begin{exercise}{310}{区域分割与积分份额}
设 $D=[0,1]^2$，以直线 $x+y=1$ 分割为 $D_-=\{(x,y)\in D:x+y\leq1\}$ 与 $D_+=D\setminus D_-$。计算 $x+y$ 在 $D,D_-,D_+$ 上的积分，并判断两个等面积部分的积分是否相等。
\end{exercise}

\begin{exercise}{311}{含参积分与零参数}
设 $\alpha\in\mathbb R$，求 $I(\alpha)=\int_0^1\int_0^1xe^{\alpha xy}\,\dd y\dd x$。分别处理 $\alpha=0$ 与 $\alpha\ne0$，求 $I(1)$，并检验所得表达式在 $\alpha=0$ 处是否连续。
\end{exercise}

\begin{exercise}{312}{三角形域与尺度参数}
设 $a>1$，$D_a$ 为直线 $x=a,y=1,y=x$ 围成的闭区域。求 $\iint_{D_a}xy\,\dd x\dd y$，给出 $a=2$ 的结果，并求 $xy$ 在 $D_a$ 上的平均值及其在 $a\to1^+$ 时的极限。
\end{exercise}

\begin{exercise}{313}{曲边域与加权平均}
设 $D$ 为 $y=\sqrt{x}$ 与 $y=x^2$ 围成的闭区域，面密度为 $\rho(x,y)=\sqrt y$。计算其质量与关于 $y$ 轴的一阶矩 $\iint_Dx\sqrt y\,\dd x\dd y$，进而求质心横坐标。须写出两曲线交点及积分限。
\end{exercise}

\begin{exercise}{314}{线性密度与形心比较}
设 $a>0$，$D_a=\{(x,y):x\geq0,y\geq0,x+y\leq a\}$，面密度 $\rho=3x+2y$。求质量 $M(a)$ 与质心横坐标，给出 $M(2)$，并比较质心横坐标与均匀薄片的形心横坐标。
\end{exercise}

\begin{exercise}{315}{三角形积分的参数表达}
设 $a>0$，$D_a=\{(x,y):0\leq y\leq x\leq a\}$。求 $I(a)=\iint_{D_a}x\cos(x+y)\,\dd x\dd y$，计算 $I(\pi)$；利用积分定义判断 $I(a)$ 在足够小的正 $a$ 下的符号，检验公式。
\end{exercise}

\begin{exercise}{316}{双曲线边界与参数}
设 $a>1$，$D_a$ 为 $x=a,xy=1,y=x$ 围成的第一象限闭区域。求 $\iint_{D_a}x^2/y^2\,\dd x\dd y$，给出 $a=2$ 的结果；说明该积分在每个固定 $a>1$ 时不存在奇点。
\end{exercise}

\begin{exercise}{317}{非均匀薄片的一阶矩}
设 $D=\{(x,y):0\leq x\leq1,x\leq y\leq5x\}$，密度为 $x+6y$。计算质量与关于 $y$ 轴的一阶矩，求质心横坐标；判断能否仅用几何形心代替质心，并用计算结果说明。
\end{exercise}

\begin{exercise}{318}{梯形区域与边界参数}
设 $0<h<\pi-1$，$D_h=\{(x,y):0\leq x\leq1,0\leq y\leq h+x\}$。求 $I(h)=\iint_{D_h}(1+x)\sin y\,\dd x\dd y$，给出 $h=1$ 的结果，并从积分区域与被积函数判断 $I(h)$ 的单调性。
\end{exercise}

\begin{exercise}{319}{缩放区域与积分符号}
设 $a>0$，$D_a=\{(x,y):0\leq y\leq a,y/2\leq x\leq y\}$。求
$I(a)=\iint_{D_a}(x^2+y^2-x)\,\dd x\dd y$，计算 $I(2)$，并完整讨论 $I(a)$ 的正负与零点。
\end{exercise}

\begin{exercise}{320}{平行四边形上的变量代换}
设 $D=\{(x,y):1\leq y\leq3,y-1\leq x\leq y\}$。以 $u=y-x,v=y$ 把积分化为矩形域积分，计算 $\iint_D(x^2+y^2)\,\dd x\dd y$ 与 $\iint_D(y^2-x^2)\,\dd x\dd y$，并说明后者为何严格为正。
\end{exercise}

\begin{exercise}{321}{三角形薄片的一阶矩}
设 $D$ 为顶点 $O(0,0),A(1,2),B(2,1)$ 的三角形闭区域。求 $\iint_Dx\,\dd x\dd y$，再利用 $x,y$ 互换的对称性求均匀薄片的形心，并用三角形顶点坐标的平均值检验结果。
\end{exercise}

\begin{exercise}{322}{积分次序与非初等原函数}
设 $D=\{(x,y):0\leq y\leq2,1\leq x\leq y+1\}$。将 $\iint_D\cos(y^2)\,\dd x\dd y$ 写成两种次序的累次积分，选取能直接计算的次序求值，并说明为何无须求 $\cos(y^2)$ 的初等原函数。
\end{exercise}

\begin{exercise}{323}{可去奇点与曲边区域}
令 $f(x)=\sin x/x$（$x\ne0$），$f(0)=1$。设 $D$ 由 $y=x,y=x^2$ 围成。计算 $\iint_Df(x)\,\dd x\dd y$，并证明在整个 $D$ 上被积函数连续；若仅改变原点处的取值，积分是否改变？
\end{exercise}

\begin{exercise}{324}{三角形域与符号估计}
设 $D=\{(x,y):0\leq x\leq y\leq1\}$。求 $I=\iint_D(e^{-y^2}-1)\,\dd x\dd y$；不依赖准确结果，直接利用被积函数证明 $-1/2<I<0$，再检验所求值。
\end{exercise}

\begin{exercise}{325}{密度积分与高斯型因子}
设 $D=\{(x,y):0\leq x\leq y\leq1\}$。求 $\iint_Dx^2e^{-y^2}\,\dd x\dd y$；若薄片密度为 $e^{-y^2}$，再求其关于 $y$ 轴的转动惯量与质量之比。
\end{exercise}

\begin{exercise}{326}{曲线间宽度与精确积分}
设 $0<a<1$，$D_a=\{(x,y):a\leq x\leq1,x^2\leq y\leq x\}$。求 $I(a)=\iint_{D_a}e^{x^2}\,\dd x\dd y$ 的一元积分表达式；在 $a=1/2$ 时将其化为含 $\int_{1/2}^1e^{x^2}\,\dd x$ 的表达式，并证明 $I'(a)<0$。允许保留所指定的非初等定积分。
\end{exercise}

\begin{exercise}{327}{换序与线性分解}
将 $I=\int_0^1\dd x\int_x^1(y+xye^{y^2})\,\dd y$ 交换积分次序并计算。另求 $y$ 项与 $xye^{y^2}$ 项对积分的贡献，比较两项大小。
\end{exercise}

\begin{exercise}{328}{参数上限与积分求导}
设 $a>0$，$I(a)=\int_0^a\dd x\int_x^a\frac{x}{\sqrt{1+y^3}}\,\dd y$。交换次序求 $I(a)$，给出 $I(1)$；由积分表达式与所得公式分别求 $I'(a)$，检验两者一致。
\end{exercise}

\begin{exercise}{329}{角点反常积分与换序}
设 $D=\{(x,y):1<x\leq2,0\leq y\leq\ln x\}$，定义
\[
 I=\lim_{\varepsilon\to0^+}\int_{1+\varepsilon}^2\dd x\int_0^{\ln x}\frac{e^{xy}}{x^x-1}\,\dd y.
\]
证明该极限存在并计算；写出形式上先对 $x$ 积分的表达式，指出角点 $(1,0)$ 使本题不能直接按连续函数的有界区域积分处理的原因。
\end{exercise}

\begin{exercise}{330}{换序与分部积分}
设 $f\in C^3[0,2]$，$f(0)=f'(0)=f''(0)=-1$，$f(2)=-1/2$。计算
\[
 I=\int_0^2\dd x\int_0^x\sqrt{(2-x)(2-y)}f'''(y)\,\dd y.
\]
先将 $I$ 化为仅含端点值 $f(0),f'(0),f''(0),f(2)$ 的表达式；说明为什么未知的 $f'(2),f''(2)$ 不影响结果。
\end{exercise}

\begin{exercise}{331}{平面支集与卷积截面}
设 $f(x)=g(x)=a$（$0\leq x\leq1$），在其余点取零，其中 $a>0$。求
$H(y)=\int_{-\infty}^{+\infty}f(x)g(y-x)\,\dd x$ 的分段表达式，进而求 $\iint_{\mathbb R^2}f(x)g(y-x)\,\dd x\dd y$。须从两个因子的非零条件确定支集。
\end{exercise}

\begin{exercise}{332}{函数支集与圆外区域}
设 $f(x,y)=x^2y$（$1\leq x\leq2,0\leq y\leq x$），其余点取零。求 $f$ 在 $D=\{x^2+y^2\geq2x\}$ 上的积分，并求在补集上的积分；用两者之和检验对整个支集的积分。
\end{exercise}

\begin{exercise}{333}{绝对值分区与最小值}
设 $0\leq c\leq1$，$D=[-1,1]\times[0,1]$。求 $I(c)=\iint_D|y-cx^2|\,\dd x\dd y$，计算 $I(1)$，并求使 $I(c)$ 最小的 $c$。写明去掉绝对值所需的分界曲线。
\end{exercise}

\begin{exercise}{334}{最大值函数与分界双曲线}
设 $0<c\leq4$，$D=[0,2]^2$。求 $I(c)=\iint_D\max\{xy,c\}\,\dd x\dd y$，给出 $I(1)$；从几何分区说明 $c=4$ 时公式应为何值，并检验结果。
\end{exercise}

\begin{exercise}{335}{对角线分割与连续参数}
设 $\alpha\in\mathbb R$，$D=[0,1]^2$。沿 $y=x$ 分区求
$I(\alpha)=\iint_De^{\alpha\max\{x^2,y^2\}}\,\dd x\dd y$，给出 $I(1)$，并用极限检验 $\alpha=0$ 时的值。
\end{exercise}

\begin{exercise}{336}{根式边界与缩放规律}
设 $R>0$，$D_R=\{(x,y):\sqrt{|x|}+\sqrt{|y|}\leq R\}$。求 $D_R$ 的面积及
$\iint_{D_R}(\sqrt{|x|}+\sqrt{|y|})\,\dd x\dd y$，给出 $R=1$ 的结果；说明两者分别按 $R$ 的几次幂缩放，并验证函数平均值小于 $R$。
\end{exercise}

\subsubsection{极坐标与圆形边界}

\begin{exercise}{337}{扇环域与角函数平均}
设 $0<\alpha<\beta<\pi/2$，$D=\{(r\cos\theta,r\sin\theta):1\leq r\leq3,\alpha\leq\theta\leq\beta\}$。
求 $\iint_D\arctan(y/x)\,\dd x\dd y$ 及该函数在 $D$ 上的平均值；给出边界直线为 $x=\sqrt3y$ 与 $y=\sqrt3x$ 时的结果。
\end{exercise}

\begin{exercise}{338}{圆域根式积分}
设 $D=\{(x,y):x^2+y^2\leq1,x\geq0,y\geq0\}$。计算
$I=\iint_D\sqrt{\frac{1-x^2-y^2}{1+x^2+y^2}}\,\dd x\dd y$；用被积函数的范围独立证明 $0<I<\pi/4$，检验准确值。写出极坐标中的面积因子。
\end{exercise}

\begin{exercise}{339}{半球体积与半径确定}
设 $a,V>0$，$D_a=\{x^2+y^2\leq a^2\}$。求
$\iint_{D_a}\sqrt{a^2-x^2-y^2}\,\dd x\dd y$ 关于 $a$ 的表达式；证明令其等于 $V$ 时半径 $a$ 唯一确定，求 $V=\pi$ 时的 $a$，并给出体积解释。
\end{exercise}

\begin{exercise}{340}{四分之一圆域上的平均距离}
设 $a>0$。计算 $\int_0^a\dd x\int_0^{\sqrt{a^2-x^2}}\sqrt{x^2+y^2}\,\dd y$，进而求四分之一圆盘上点到原点距离的平均值，说明所得平均值与 $a$ 的比例为何不随 $a$ 改变。
\end{exercise}

\begin{exercise}{341}{平移圆域与积分零点}
设 $h\in\mathbb R$，$D_h=\{(x,y):x^2+(y-h)^2\leq1\}$。求 $I(h)=\iint_{D_h}(x^2-y)\,\dd x\dd y$，给出 $h=0$ 的结果，并求所有使 $I(h)=0$ 的 $h$。
\end{exercise}

\begin{exercise}{342}{环域与径向对数积分}
设 $R>1$，$D_R=\{(x,y):1\leq x^2+y^2\leq R^2\}$。求 $I(R)=\iint_{D_R}x^2\ln(x^2+y^2)\,\dd x\dd y$，给出 $R=e$ 的结果；利用对称性说明把 $x^2$ 换为 $y^2$ 后积分不变，并求换为 $x^2+y^2$ 后的值。
\end{exercise}

\begin{exercise}{343}{互换对称性与孤立不连续点}
设 $D=\{x^2+y^2\leq1,x\geq0,y\geq0\}$。将
$x\cos(x^2+y^2)/(x+y)$ 在原点定义为零，求其在 $D$ 上的积分；判断此延拓在原点是否连续，说明该点是否影响积分。可将分子 $x$ 与 $y$ 互换后相加简化计算。
\end{exercise}

\begin{exercise}{344}{原点奇性与径向截断}
设 $D=\{(x,y):0\leq x\leq1,0\leq y\leq x\}$。用截断区域 $D\cap\{x^2+y^2\geq\varepsilon^2\}$ 定义并计算
$\iint_D(x^2+y^2)^{-1/2}\,\dd x\dd y$，证明极限有限，并说明极坐标面积因子如何抵消原点奇性。
\end{exercise}

\begin{exercise}{345}{扇形边界下的径向矩}
设 $a>0$，$D_a=\{(x,y):0\leq x\leq a,0\leq y\leq x\}$。计算 $\iint_{D_a}\sqrt{x^2+y^2}\,\dd x\dd y$，求 $a=1$ 的结果；求距离的平均值，并判断其是否小于区域内最大距离。
\end{exercise}

\begin{exercise}{346}{参数扇域与退化积分}
设 $a\geq1$，定义 $I(a)=\int_0^1\dd x\int_x^{ax}\frac1{\sqrt{x^2+y^2}}\,\dd y$，原点处按反常积分极限理解。用适当换元求 $I(a)$，证明其收敛并讨论单调性；求 $a=1$ 的值，解释积分区域退化时结果如何理解。
\end{exercise}

\begin{exercise}{347}{两圆之间的极坐标区域}
设 $D$ 为第一象限内由 $y=0,y=x$ 及圆弧 $x^2+y^2=4x,x^2+y^2=9x$ 围成的闭区域。求其极坐标范围，计算 $\iint_Dy^2/x^2\,\dd x\dd y$；另求 $D$ 的面积，从而求 $(y/x)^2$ 的平均值。
\end{exercise}

\begin{exercise}{348}{平移旋转与半圆积分}
设 $D=\{(x,y):(x-1)^2+(y-1)^2\leq2,y\geq x\}$。作适当的平移与旋转，把 $x-y$ 化为一个坐标，求 $\iint_D(x-y)\,\dd x\dd y$，并求 $x-y$ 在 $D$ 上的平均值。
\end{exercise}

\begin{exercise}{349}{偏心圆盘的矩}
设 $a>0$，$D_a=\{(x,y):x^2+y^2\leq ax\}$。求 $\iint_{D_a}(x+y^2)\,\dd x\dd y$，给出 $a=1$ 的结果；分别用平移圆心后的对称性确定两个分项的缩放次数。
\end{exercise}

\begin{exercise}{350}{圆域积分方程的可解性}
设 $D=\{(x,y):x^2+y^2\leq y,x\geq0\}$，$\lambda\in\mathbb R$。讨论连续函数方程
\[
 f(x,y)=\sqrt{1-x^2-y^2}-\lambda\iint_Df(u,v)\,\dd u\dd v
\]
的存在性与唯一性，存在时求其解；给出 $\lambda=8/\pi$ 的结果。必须算出区域面积及根式项的积分。
\end{exercise}

\begin{exercise}{351}{最小值函数与径向分区}
设 $D=\{x^2+y^2\leq3/16\}$。求两函数 $\sqrt{3/16-x^2-y^2}$ 与 $2(x^2+y^2)$ 在 $D$ 内的相等点集，计算其最小值的二重积分；再利用“最小值加最大值等于两者之和”求最大值的积分。
\end{exercise}

\begin{exercise}{352}{绝对值与挖去圆盘}
设 $D=\{(x,y):x^2+y^2\geq x,|x|\leq1,|y|\leq1\}$。把 $D$ 解释为正方形挖去一个圆盘后的区域，求 $\iint_D|xy|\,\dd x\dd y$；同时求 $\iint_Dxy\,\dd x\dd y$，说明两个积分能否使用相同的奇偶性结论。
\end{exercise}

\subsubsection{一般换元与无界积分}

\begin{exercise}{353}{线性变换与雅可比}
设 $a,b>0$，$D_{a,b}=\{0\leq x+y\leq a,0\leq x-y\leq b\}$。求
$\iint_{D_{a,b}}(x+y)\sin(x-y)\,\dd x\dd y$，给出 $a=b=\pi$ 的结果。写出正、逆变换及雅可比绝对值，并解释为何不能遗漏该因子。
\end{exercise}

\begin{exercise}{354}{椭圆变换与幂参数}
设 $a,b>0$，$p\geq0$，$D=\{x^2/a^2+y^2/b^2\leq1\}$。求
$\iint_D(x^2/a^2+y^2/b^2)^p\,\dd x\dd y$，分别给出 $p=0,1$ 的结果，并从中求原二次函数在椭圆内的平均值。
\end{exercise}

\begin{exercise}{355}{和与比值变量的变换}
设 $D=\{x\geq0,y\geq0,x+y\leq1\}$。在原点给 $e^{x/(x+y)}$ 任意有限值。用 $u=x+y,v=x/(x+y)$ 计算积分，并求以该函数为面密度时的总质量及 $x+y$ 的加权平均值；说明变换在原点失效是否影响结果。
\end{exercise}

\begin{exercise}{356}{幂乘积变换与区域面积}
设 $D$ 为第一象限内由 $xy=4,xy=8,xy^3=5,xy^3=15$ 围成的闭区域。取 $u=xy,v=xy^3$，证明逆变换在相应矩形内唯一，求雅可比绝对值并计算 $D$ 的面积；再求 $\iint_Dxy\,\dd x\dd y$。
\end{exercise}

\begin{exercise}{357}{菱形域的积分恒等式}
设 $f$ 在 $[-1,1]$ 上连续，$D=\{|x|+|y|\leq1\}$。证明
$\iint_Df(x+y)\,\dd x\dd y=\int_{-1}^1f(u)\,\dd u$。
进一步设 $g$ 也连续，推导 $\iint_Df(x+y)g(x-y)\,\dd x\dd y$ 的乘积公式，并检验 $f=g\equiv1$ 时的面积。
\end{exercise}

\begin{exercise}{358}{无界曲边区域的收敛性}
设 $D=\{(x,y):x\geq0,4x^2\leq y\leq9x^2\}$。先以 $0\leq y\leq R$ 截断，计算 $\iint_Dxe^{-y^2}\,\dd x\dd y$ 并证明其收敛；再判断将 $e^{-y^2}$ 换为 $1$ 后积分是否收敛。
\end{exercise}

\begin{exercise}{359}{无界扇形上的高斯积分}
设 $\alpha>0$。证明并计算
$I(\alpha)=\int_0^{+\infty}\dd x\int_x^{\sqrt3x}e^{-\alpha(x^2+y^2)}\,\dd y$，给出 $I(1)$；利用平面缩放独立说明 $I(\alpha)$ 与 $\alpha$ 的关系。
\end{exercise}

\begin{exercise}{360}{第一象限反常积分的参数界限}
设 $p\in\mathbb R$，以半径 $R$ 的第一象限圆盘递增截断定义
$I(p)=\iint_{x,y\geq0}(1+x^2+y^2)^{-p}\,\dd x\dd y$。
求积分收敛的全部 $p$ 及收敛时的值；特别证明 $I(2)=\pi/4$，并分别说明临界值与临界值以下的发散情形。
\end{exercise}

\begin{exercise}{361}{换元区域与雅可比的双重核查}
设 $D=\{0\leq y\leq x\leq1\}$，取 $u=x+y,v=x-y$。有人将 $D$ 的像写成矩形 $[0,2]\times[0,1]$，并把雅可比因子取为 $1$。求真实像区域、逆变换与正确的面积因子，计算 $\iint_D(x+y)\,\dd x\dd y$；再用原坐标积分检验，分别说明两项错误各自改变了什么。
\end{exercise}

\begin{exercise}{362}{圆域奇点的两个临界指数}
设 $p\in\mathbb R$，$D=\{x^2+y^2\leq1\}$。以 $D\cap\{x^2+y^2\geq\varepsilon^2\}$ 截断，分别讨论
\[
 \iint_D(x^2+y^2)^{-p/2}\,\dd\sigma,\qquad
 \iint_D\frac{x^2}{(x^2+y^2)^{p/2}}\,\dd\sigma
\]
的收敛参数范围及收敛时的值；写出临界参数处的截断表达式，解释分子 $x^2$ 如何改变奇性阶数。
\end{exercise}

\subsection{三重积分}

\begin{exercise}{363}{空间区域与幂函数积分}
设 $\Omega=\{(x,y,z):0\leq y\leq x\leq1,\ 0\leq z\leq xy\}$。
写出先对 $z$ 积分与先对 $y$ 积分的两种累次积分，计算
$\iiint_\Omega xy^2z^3\,\dd x\dd y\dd z$，并说明曲面 $z=xy$ 与 $z=0$ 如何确定区域的下边界 $y=0$。
\end{exercise}

\begin{exercise}{364}{三重积分的次序选择}
设 $\Omega=\{(x,y,z):0\leq x\leq\pi/2,\ 0\leq y\leq\sqrt{x},\ 0\leq z\leq\pi/2-x\}$。
将 $\sin x/x$ 在 $x=0$ 处作连续延拓，计算
\[
 \iiint_\Omega\frac{y\sin x}{x}\,\dd x\dd y\dd z,
\]
并解释先对 $y,z$ 积分为何能消去分母 $x$ 带来的计算困难。
\end{exercise}

\begin{exercise}{365}{椭球截面与参数积分}
设 $0<h\leq3$，$\Omega_h=\{(x,y,z):x^2+y^2/4+z^2/9\leq1,\ 0\leq z\leq h\}$。
用水平截面法求 $J(h)=\iiint_{\Omega_h}z^2\,\dd V$，求 $J(1)$，并说明 $J'(h)$ 的几何意义。
\end{exercise}

\begin{exercise}{366}{球面与抛物面的交域}
设 $\Omega$ 由 $z=\sqrt{2-x^2-y^2}$ 与 $z=x^2+y^2$ 围成。求两曲面的交线及区域的柱面坐标范围，计算体积与 $\iiint_\Omega z\,\dd V$，进而求均匀物体的形心高度。
\end{exercise}

\begin{exercise}{367}{抛物面内水平截层}
设 $0<a<b$，$\Omega_{a,b}=\{(x,y,z):a\leq z\leq b,x^2+y^2\leq z\}$。求 $\iiint_{\Omega_{a,b}}(x^2+y^2)^2\,\dd V$，给出 $a=1,b=2$ 的结果；求该函数在区域上的平均值。
\end{exercise}

\begin{exercise}{368}{旋转曲面与参数截面}
将平面 $x=0$ 内的抛物线 $y^2=2z$ 绕 $z$ 轴旋转，所得曲面与平面 $z=H$（$H>0$）围成 $\Omega_H$。求曲面方程与 $\iiint_{\Omega_H}(x^2+y^2+z^2)\,\dd V$，并给出 $H=1$ 的结果。
\end{exercise}

\begin{exercise}{369}{球域幂函数与坐标面}
设 $R>0$，$\Omega_R=\{x^2+y^2+z^2\leq R^2,x,y,z\geq0\}$。计算 $\iiint_{\Omega_R}xyz\,\dd V$，给出 $R=1$ 的结果；若改在整个球域上积分，结果如何？说明区域改变对奇偶性的影响。
\end{exercise}

\begin{exercise}{370}{偏心球与圆锥的交域}
设 $\Omega=\{x^2+y^2+z^2\leq2z,z\geq\sqrt{x^2+y^2}\}$。写出球面坐标范围，计算 $\iiint_\Omega z^3\,\dd V$；再用水平截面说明截面半径在何处改变表达式，写出另一种计算式。
\end{exercise}

\begin{exercise}{371}{半球与卦限的对称性辨析}
设 $R>0$，$\Omega_1=\{x^2+y^2+z^2\leq R^2,z\geq0\}$，$\Omega_2=\Omega_1\cap\{x\geq0,y\geq0\}$。对 $f=x,y,z,xyz$，分别判断
\[
 \iiint_{\Omega_1}f\,\dd V=4\iiint_{\Omega_2}f\,\dd V
\]
是否成立，给出理由；对成立的非零情形，计算两边的共同值。
\end{exercise}

\begin{exercise}{372}{球域的二次型积分}
设 $a,b,c>0$，$\Omega=\{x^2+y^2+z^2\leq1\}$。利用坐标互换求
$\iiint_\Omega(x^2/a^2+y^2/b^2+z^2/c^2)\,\dd V$；进一步求加入混合项 $\lambda xy+\mu yz+\nu zx$ 后的积分，说明这些系数为何不影响结果。
\end{exercise}

\begin{exercise}{373}{球内二次矩与转动惯量}
设 $R>0$，$\Omega=\{x^2+y^2+z^2\leq R^2\}$。求 $\iiint_\Omega y^2\,\dd V$，并由对称性求密度为 $1$ 的实心球关于任一直径的转动惯量；检验结果关于 $R$ 的缩放次数。
\end{exercise}

\begin{exercise}{374}{圆锥截球与平均高度}
设 $\Omega$ 由圆锥面 $z=\sqrt{x^2+y^2}$ 与上半球面 $z=\sqrt{1-x^2-y^2}$ 围成。先说明 $x$ 项的积分是否为零，再求 $\iiint_\Omega(x+z)\,\dd V$ 与该区域的平均高度。
\end{exercise}

\begin{exercise}{375}{径向积分比值的单调性}
设 $f$ 在 $[0,+\infty)$ 上连续且严格为正，$t>0$。令
\[
 F(t)=\frac{\iiint_{x^2+y^2+z^2\leq t^2}f(x^2+y^2+z^2)\,\dd V}
 {\iint_{x^2+y^2\leq t^2}f(x^2+y^2)\,\dd\sigma}.
\]
把 $F(t)$ 化为一元积分之比，证明其在 $(0,+\infty)$ 上严格递增，并求 $\lim_{t\to0^+}F(t)/t$。证明中须明确用到 $f>0$ 的位置。
\end{exercise}

\begin{exercise}{376}{小圆柱上的积分极限}
设 $f$ 在 $0$ 的右邻域连续可微，$f(0)=0$，$H>0$，$\Omega_t=\{x^2+y^2\leq t^2,0\leq z\leq H\}$。求
$\lim_{t\to0^+}t^{-4}\iiint_{\Omega_t}f(x^2+y^2)\,\dd V$，给出 $H=1$ 的结果；用 $f(u)=f'(0)u+o(u)$ 的余项估计证明极限。
\end{exercise}

\begin{exercise}{377}{球冠的坐标范围与截面验证}
设 $0<h<1$，$\Omega_h=\{x^2+y^2+z^2\leq1,z\geq h\}$。分别用水平截面与球面坐标写出体积积分，明确球面坐标中径向变量的下界及极角范围。计算体积与 $\iiint_{\Omega_h}z\,\dd V$，由两种表达式检验结果，并判断 $h\to1^-$ 时平均高度的极限。
\end{exercise}

\begin{exercise}{378}{非均匀抛物体的参数质心}
设 $\Omega=\{x^2+y^2\leq z\leq1\}$，体密度为 $\rho_\alpha=1+\alpha z$。
先求密度在整个区域上非负的全部 $\alpha$。在该范围内求质量与全部质心坐标，证明质心高度随 $\alpha$ 严格增加，求其最小值与 $\alpha\to+\infty$ 时的极限，并确认质量不会为零。
\end{exercise}

\subsection{重积分的应用}

\begin{exercise}{379}{坐标面截三角形的面积}
设 $a,b,c>0$。求平面 $x/a+y/b+z/c=1$ 被三个坐标面截出的三角形面积。分别写出向 $xOy$ 面投影与向 $yOz$ 面投影的面积积分，并验证两种结果相同。
\end{exercise}

\begin{exercise}{380}{柱内球面的面积}
设 $a>0$。求球面 $x^2+y^2+z^2=a^2$ 满足 $x^2+y^2\leq ax$ 部分的面积，明确是否包含上、下半球的部分；再求该面积占整个球面面积的比例，并说明比例为何与 $a$ 无关。
\end{exercise}

\begin{exercise}{381}{球带面积与球面总面积}
设 $R>0$，$-R\leq a<b\leq R$。用曲面面积积分求球面 $x^2+y^2+z^2=R^2$ 上 $a\leq z\leq b$ 的球带面积，进而推导整球表面积；若积分表达式在极点或赤道出现无穷因子，说明相应极限的处理方法。
\end{exercise}

\begin{exercise}{382}{曲边薄片的质量与质心}
设 $D=\{(x,y):0\leq x\leq1,x^2\leq y\leq x\}$，面密度为 $\mu(x,y)=x^2y$。求质量及质心坐标；再求同一区域均匀薄片的形心，比较两个横坐标，并从密度随 $x$ 的变化解释结果。
\end{exercise}

\begin{exercise}{383}{偏心球体的非均匀质心}
设球体 $\Omega=\{x^2+y^2+z^2\leq z\}$ 的体密度为 $\rho=x^2+y^2+z^2$。利用对称性确定质心的两个坐标，计算质量及质心竖坐标，并比较它与均匀球体的球心高度。
\end{exercise}

\begin{exercise}{384}{抛物体截面的形心}
设 $H>0$，$\Omega_H=\{(x,y,z):x^2+y^2\leq z\leq H\}$。求其体积与全部形心坐标，给出 $H=1$ 的结果；解释形心高度为何不能直接取上下端高度的算术平均值。
\end{exercise}

\begin{exercise}{385}{斜圆锥的截面与形心}
设 $\Omega=\{(x,y,z):0\leq z\leq1,x^2+(y-z)^2\leq(1-z)^2\}$ 为均匀锥体。求水平截面的面积及圆心，利用截面法求全部形心坐标，并证明形心位于底面圆心与顶点的连线上。
\end{exercise}

\begin{exercise}{386}{椭圆薄片的惯量}
设 $a,b>0$，$D=\{x^2/a^2+y^2/b^2\leq1\}$，面密度为 $1$。求薄片关于 $y$ 轴、$x$ 轴的转动惯量与关于原点的极惯量；比较前两者并给出相等的充要条件。
\end{exercise}

\begin{exercise}{387}{平行直线的惯量最小值}
设 $D=\{(x,y):-1\leq x\leq1,x^2\leq y\leq1\}$，面密度为常数 $\mu>0$。求关于直线 $y=c$ 的转动惯量 $I(c)$，给出 $c=-1$ 的结果，并求使 $I(c)$ 最小的 $c$，说明其与质心的关系。
\end{exercise}

\begin{exercise}{388}{抛物面下立体的轴惯量}
设 $a,\rho>0$，$\Omega=\{(x,y,z):|x|\leq a,|y|\leq a,0\leq z\leq x^2+y^2\}$，体密度为 $\rho$。求质量及关于 $z$ 轴的转动惯量，给出惯量与质量之比，并检验其量纲。
\end{exercise}

\begin{exercise}{389}{圆柱体的轴向引力}
设 $R,h>0$，均匀圆柱体 $\Omega=\{x^2+y^2\leq R^2,0\leq z\leq h\}$ 的体密度为 $1$，引力常数为 $G$。求其对 $M_0(0,0,a)$（$a>h$）处单位质量质点的引力向量；先用对称性确定方向，再计算非零分量，并检验 $a\to+\infty$ 时引力大小与 $G\pi R^2h/a^2$ 的比值。
\end{exercise}

\begin{exercise}{390}{曲面面积的投影因子与边界极限}
设 $S:z=\sqrt{1-x^2-y^2}$ 为单位上半球面。解释为何投影圆盘面积不能直接作为 $S$ 的面积。
先在 $x^2+y^2\leq(1-\varepsilon)^2$（$0<\varepsilon<1$）上计算对应球面面积，再令 $\varepsilon\to0^+$；指出面积因子在赤道附近虽无界，面积仍有限的具体积分依据。
\end{exercise}

\begin{exercise}{391}{椭圆薄片的非对称密度与中心惯量}
设 $a,b>0$，$D=\{x^2/a^2+y^2/b^2\leq1\}$，面密度为 $1+\lambda x/a$。
求使密度非负的全部 $\lambda$，在该范围内求质量、质心、关于 $y$ 轴的惯量，以及关于过质心且平行于 $y$ 轴的直线的惯量。证明最后一个惯量的最大值与最小值分别在何处取得，并解释质量为何不随 $\lambda$ 改变。
\end{exercise}
